% -----------------------------------------------------------------------------
%  entries/q027.tex
%  Issued 22/10/2026. Second entry of the GKP / magic / hypergraph sequence;
%  grew out of Q025(d)(iii). All items solved before issue (pure algebra; the
%  transformation rule of (a) checked on the quarter turn F of Q025).
% -----------------------------------------------------------------------------
\begin{entry}{Q027}{No Gaussian gate goes beyond Clifford}

\entrymeta{
  topic   = {GKP code; Gaussian unitaries; symplectic group; logical Clifford group; universality},
  est     = {$\sim$2 hours including the write-up: (a) $\sim$45 min, (b) $\sim$25 min, (c) $\sim$50 min},
  origin  = Original,
  source  = {Original problem. Gaussian unitaries as affine symplectic maps:
    C.~Weedbrook \emph{et al.}, \emph{Gaussian quantum information},
    Rev.~Mod.~Phys.\ \textbf{84}, 621 (2012), Sec.~II.A.2},
  url     = https://arxiv.org/abs/1110.3234,
  prereq  = {\qref{Q025}, in particular \eqref{q025:shifts}, \eqref{q025:weyl} and
    \eqref{q025:code}},
  status  = drafted,
}

\begin{problem}
\emph{Context.} \qref{Q025} showed that the logical Fourier, sum and phase gates of the
GKP code are Gaussian unitaries. This entry proves the converse: a Gaussian unitary that
preserves the code can only ever act as a logical Clifford gate. The Clifford group is
finite up to phases, so Gaussian optics, however well controlled, cannot by itself compute
universally on GKP qudits. The missing non-Clifford gates must draw on a non-Gaussian
resource, which is where the hypergraph interactions of \qref{Q022} and magic states enter
in the following entries. Along the way you will find which symplectic maps of phase space
are compatible with the code lattice.

\emph{Setting.} Use the conventions of \qref{Q025}: one mode, the shifts $X(a)$ and $Z(b)$,
$\ell=\sqrt{2\pi/d}$, the stabilizers $S_q,S_p$ and the logical Paulis $\bar X,\bar Z$. For
$\xi=(\xi_q,\xi_p)\in\RR^2$ write
\begin{equation}
  D(\xi):=X(\xi_q)\,Z(\xi_p),\qquad D(\xi)^\dagger\,\hat r\,D(\xi)=\hat r+\xi,\qquad
  \hat r:=\begin{pmatrix}\hat q\\ \hat p\end{pmatrix},
  \label{q026:shift}
\end{equation}
so that $D(\xi)$ shifts phase space by $\xi$. A unitary $G$ is \emph{Gaussian} if it is
generated by Hamiltonians at most quadratic in $\hat q$ and $\hat p$. Equivalently (Weedbrook
\emph{et al.}, Sec.~II.A.2), it acts on the quadratures as an affine map,
\begin{equation}
  G^\dagger\hat r\,G=S\hat r+c,\qquad
  S\ \text{a real }2\times2\text{ matrix with }\det S=1,\quad c\in\RR^{2};
  \label{q026:gaussian}
\end{equation}
for one mode, $S$ is symplectic exactly when $\det S=1$. You may use that a unitary $U$ with
$U^\dagger\hat qU=\hat q$ and $U^\dagger\hat pU=\hat p$ is a multiple of the identity: the
canonical commutation relations are represented irreducibly on $L^2(\RR)$ (B.~C.~Hall,
\emph{Quantum Theory for Mathematicians}, Springer 2013, the chapter on the Stone--von
Neumann theorem). A unitary on the code space is a \emph{logical Clifford gate} if
conjugation by it maps every logical Pauli operator $\bar X^{m}\bar Z^{m'}$ to a logical
Pauli operator, up to a phase.

\begin{enumerate}[label=\textbf{(\alph*)},leftmargin=*]
  \item \emph{Gaussian unitaries move shifts.} Show that for every Gaussian unitary $G$
        and every $\xi\in\RR^{2}$,
        \begin{equation}
          G^\dagger D(\xi)\,G=e^{i\theta}\,D\bigl(S^{-1}\xi\bigr)
          \label{q026:conjugation}
        \end{equation}
        for some phase $\theta$.

  \item \emph{Which shifts preserve the code.} Show that $D(\xi)$ maps the code space into
        itself if and only if $\xi\in\ell\ZZ^{2}$, that is, if and only if $D(\xi)$ is, up to
        a phase, a logical Pauli operator $\bar X^{m}\bar Z^{m'}$.

  \item \emph{The lattice condition.} Let $G$ be a Gaussian unitary that maps the code
        space onto itself.
        \begin{enumerate}[label=(\roman*)]
          \item Show that $S$ has integer entries, that is, $S\in SL(2,\ZZ)$, and that $G$
                acts on the code as a logical Clifford gate, sending the labels $(m,m')$ of
                a logical Pauli operator to $S(m,m')$ modulo $d$. Check the condition for $F$ the quarter turn from \qref{Q025}
          \item Read off the integer matrices of the quarter turn $F$ and of
                $e^{i\hat q^{2}/2}$ from \qref{Q025}. Is
                $S\in SL(2,\ZZ)$ also sufficient for $G$ to preserve the code? And decide whether a single-mode
                squeezer, $S=\mathrm{diag}(s,1/s)$ with $s>0$, can preserve the code.
        \end{enumerate}

\end{enumerate}
\end{problem}

\begin{solution}[Part (a)]
Let's see how this transformed operator acts, $G^\dagger\,D(\xi) G$, on $\hat{r}$ 
\begin{equation}
   G^\dagger D(\xi)^\dagger\, G\hat r G^\dagger\,D(\xi) G = 
\end{equation}
Note that the order of the action of Gaussians here is different than \eqref{q026:gaussian}. First, notice that multiplying $G^\dagger$ from the right and $G$ from the left results in
\begin{equation}
    \hat r = S G\hat r G^\dagger +c \implies G\hat r G^\dagger = S^{-1} (\hat r - c).
\end{equation}
Thus, going back to our equation 
\begin{equation}
\begin{split}
    \hat r \rightarrow& S^{-1} G^\dagger D(\xi)^\dagger\,  \hat r \,D(\xi) G -  S^{-1}c \\
    &=  S^{-1} G^\dagger   (\hat r +\xi) G -  S^{-1}c\\
    &=  S^{-1} \left(S\hat r+c\right) +S^{-1}\xi -  S^{-1}c\\
    &=  \hat r +S^{-1}\xi 
\end{split} 
\end{equation}
Which is precisely how any $e^{i\theta}\,D\bigl(S^{-1}\xi\bigr)$ acts on $\hat r$ by conjugation for any $\theta$. 

To see there is a $\theta$ that satisfies \eqref{q026:conjugation}, take $A=G^\dagger D(\xi)G$ and $B=D(S^{-1}\xi)$, from our previous calculation $A^\dagger\hat r A = B^\dagger \hat r B$. Now consider how $U=AB^\dagger$ acts by conjugation on $\hat r$: 
$$ U^\dagger\hat r U = B(A^\dagger \hat r A) B^\dagger = B(B^\dagger \hat r B) B^\dagger = \hat r
$$

And as given in the question, this means $U$ must be a multiple of the identity, so $U=e^{i\theta} \id$ for some $\theta$.

\end{solution}

\begin{solution}[Part (b)]
We can see how $D(\xi)$ acts by conjugation on $S_q,S_p$, using the relation \eqref{q025:weyl}
\begin{equation}
\begin{split}
    D(\xi)^\dagger S_q D(\xi) &= Z(-\xi_p)X(-\xi_q) X(d\ell)X(\xi_q) Z(\xi_p)\\
    &=   Z(-\xi_p) X(d\ell) Z(\xi_p)\\
    &=  e^{-i\xi_p d\ell} S_q
    \label{q026: conta 1}
\end{split}
\end{equation}
similarly
\begin{equation}
    D(\xi)^\dagger S_p D(\xi) =  e^{i\xi_q d\ell} S_p
    \label{q026: conta 2}
\end{equation}

Preserving the code space is equivalent to acting trivially by conjugation on the stabilizers on the code space. 
$$V\ \text{preserves the code}\iff (V^\dagger S_qV)\psi=\psi\ \text{and}\ (V^\dagger S_pV)\psi=\psi\quad\text{for every code state }\psi.$$
But the shift acts only by adding a phase by \eqref{q026: conta 1} and \eqref{q026: conta 2}, so it is necessary to fix the operator itself in this case.

Thus, we need $\xi_q d\ell = 2\pi m$, $\xi_p d\ell=2\pi m'$, for some integers $m,m'$. From the expression of $\ell$, we see that $2\pi/d=\ell^2$, such that the conditions to code preservation can be written as 
\begin{equation}
    \xi_q = \ell m,\quad \xi_p = \ell m'
\end{equation} 

\end{solution}

\begin{solution}[Part (c)]
i) First notice that $S_q = X(d\ell)= X(d\ell)Z(0\cdot d\ell)=D(d\ell e_1)$, where $e_1=(1,0)$. So we have 
\begin{equation}
    G^\dagger S_q G = G^\dagger D(d\ell\, e_1) G = e^{i\theta} D(d\ell\, S^{-1}e_1).
\end{equation}

Take $\eta = d\ell \,S^{-1}e_1$. As argued in item b), $G^\dagger S_q G$ must also stabilize the code, such that for any state in the code space $\ket{\psi}$, we must have $D(\eta)\ket{\psi} = e^{-i\theta}\ket{\psi}$. So $D(\eta)$ must not only preserve the code space, but it must act as a multiple of the identity in the code space. 

From $b$, we know that we need $\eta \in \ell \ZZ^2$ if and only if $D(\eta)$ preserves the code state. In such a case, it acts as a logical Pauli $\bar X^m \bar Z^{m'}$. Thus, it only acts as the identity when $m,m'$ are multiples of $d$, i.e., $\eta \in d\ell \ZZ^2$. This implies that $S^{-1}e_1 \in \ZZ^2$.

This demonstrates that the first column of $S^{-1}$ has integer entries. To see that the same holds for the second column, we use the same argument for $ G^\dagger S_p G$ to conclude that $S^{-1}e_2 \in \ZZ^2$. Finally, we have that $S^{-1}$ has integer entries.

As given in the question, $\det(S)=1$, so $\det(S^{-1})=1$, and then $S_{i,j}=(-1)^{i+j}(S^{-1})_{j+1,i+1}$ (being the indices summed mod 2). Therefore, $S$ has integer entries 

To see it act as Clifford gates on the code space, we must see if it maps logical Pauli gates $\bar X^m\bar Z^{m'}= D(\ell (m,m'))$ to other logical Pauli $\bar X^n,\bar Z^{n'}= D(\ell (n,n'))$. Which is the case, since from item a) 
$$G^\dagger D(\ell (m,m'))G = e^{i\theta } D(\ell\, S^{-1}(m,m'))$$, and as demonstrated $S^{-1}(m,m')\in \ZZ^2$.

ii) For the quarter turn, the corresponding $S$ matrix is
\begin{equation}
    S_F = \begin{pmatrix}
        0& -1\\
        1 & 0
    \end{pmatrix}
\end{equation}
Meanwhile, 
\begin{equation}
   e^{-i\hat q^2/2}\, \hat p\, e^{i\hat q^2/2} = \hat p + e^{-i\hat q^2/2}[\hat p\ ,\ e^{i\hat q^2/2}] = \hat p + \hat q
\end{equation}
(using that commutation by $\hat p$ act as $-i\partial_q$ for a function of $\hat q$)
And clearly $e^{-i\hat q^2/2}\, \hat q\, e^{i\hat q^2/2} = \hat q$. Thus, 
\begin{equation}
    S_{q^2} = \begin{pmatrix}
        1& 0\\
        1& 1
    \end{pmatrix}
\end{equation}

In both cases, the matrices have integer elements with determinant equal to $1$. For the quadratic phase operator $e^{i\hat q^2/2}$, as seen in \qref{q025}, it does not preserve the code for odd dimension $d$. So the condition is not sufficient. 

For the squeezing, we see that the matrix cannot have all entries being integers (for $s\neq1$), so it does not preserve the code, because despite the condition not being sufficient, it is necessary.



\end{solution}

% \begin{remark}
% \end{remark}

\begin{takeaway}
\end{takeaway}

%  Everything below is optional.  An entry with no attempt block is complete.
%  Record only what would change what you do next time: typos and arithmetic
%  slips you caught yourself are artefacts of the medium, not facts about the
%  mathematics.  Add "public," to show this block in the public build.
% \attempt{
%   date       = ,               % ISO, YYYY-MM-DD -- required if the block is used
%   time       = ,
%   difficulty = ,               % 1-5, perceived
%   howitwent  = ,
%   insight    = ,
%   note       = ,
% }

\end{entry}